% ============================================================================
% Fichier  : partie2/P02_C09_conclure_observation.tex
% Titre    : Conclure : Conformité, Incertitude et Observation
% Auteur   : MINKA MI NGUIDJOÏ Thierry Emmanuel
% Version  : 1.0 -- octobre 2026
% ============================================================================
% Sources : article A1 (HAL hal-05701270) : métriques H_q, U_q, mu_q,
% divergence de Kullback-Leibler entre cycles, quadrant diagnostique
% conformité x incertitude (quatre régions) et cas limites (ignorance
% unanime : H = 0, U = 1) ; article A0 : non-conformité connue et inconnue.
% Seuils du quadrant (mu 0,5 ; U 0,20) et conversion en niveaux de maturité :
% conventions pédagogiques du cours, déclarées comme telles. Observation en
% sept éléments : cours Audit_SI de l'auteur ; constat en cinq éléments :
% cours GCI (chapitre 11).
% ============================================================================

\chapter[Conclure]{Conclure : Conformité, Incertitude et Observation}
\label{chap:conclure}
\chaptermark{Conclure et rédiger}

\epigraph{%
    « Conclure, ce n'est pas compter les \emph{oui}.\\
    C'est savoir ce qu'ils valent. »%
}{Principe de la méthode}

\niveauChapitre{Intermédiaire}{%
    Maîtrise du Chapitre~\ref{chap:preuve} (instrument \OUI{}~/ \NON{}~/
    \NEK{}, preuves, papiers de travail)%
}

\begin{boxobjectifs}
\begin{itemize}
    \item Calculer, pour une question et pour un domaine, le taux de
          conformité~$\mu$, le taux d'incertitude~$U$ et l'entropie~$H$.
    \item Situer un domaine dans le quadrant conformité × incertitude et en
          déduire l'action appropriée.
    \item Appliquer la règle de décision du cours, qui combine toujours
          $\mu$, $U$ et $H$.
    \item Comparer deux cycles d'audit par la divergence de Kullback-Leibler.
    \item Convertir un taux de conformité en niveau de maturité, et le
          distinguer de la capacité COBIT.
    \item Rédiger une observation en sept éléments, construire le rapport et
          suivre les recommandations.
\end{itemize}
\end{boxobjectifs}

\bigskip

% ============================================================================
\section*{Question directrice : que peut-on affirmer~?}
% ============================================================================

Le terrain est terminé. AUDIX CONSEIL dispose, pour les cinq domaines de sa
mission chez LiZbiZ, de centaines de réponses et de dizaines de tests. Le
directeur général attend une phrase~: « Mes contrôles fonctionnent-ils~? »
M.~BEYALA sait qu'il doit répondre domaine par domaine, en disant à la fois
ce qui va mal et ce que l'on ne sait pas.

Ce chapitre apprend à conclure. Il définit les trois mesures
(section~\ref{sec:c09-mesures}), les lit dans un quadrant
(section~\ref{sec:c09-quadrant}), en tire une règle de décision
(section~\ref{sec:c09-regle}), compare deux cycles
(section~\ref{sec:c09-cycles}), traite de la maturité
(section~\ref{sec:c09-maturite}), puis de l'observation, du rapport et du
suivi (sections~\ref{sec:c09-observation} et~\ref{sec:c09-rapport}).

\separateur

% ============================================================================
\section{Trois Mesures}
\label{sec:c09-mesures}
% ============================================================================

\subsection{Pour une question}

Pour une question~$q$ à laquelle plusieurs répondants ont répondu, on note
$p_q(\OUI)$, $p_q(\NON)$ et $p_q(\NEK)$ les parts de chaque réponse. L'article
A1 définit trois mesures~\cite{minka2026nek}.

\begin{boxdef}
\termecle{Taux de conformité}\index{taux de conformité}~: $\mu_q =
p_q(\OUI)$, la part des réponses \OUI{}.

\medskip
\termecle{Taux d'incertitude}\index{taux d'incertitude}~: $U_q = p_q(\NEK)$,
la part des réponses \NEK{}.

\medskip
\termecle{Entropie}\index{entropie}~: dispersion des réponses, en bits,
\[
    H_q = -\sum_{m \in \{\OUI, \NON, \NEK\}} p_q(m) \log_2 p_q(m),
    \qquad 0 \leq H_q \leq \log_2 3 \approx 1{,}585,
\]
avec la convention $0 \log_2 0 = 0$.
\end{boxdef}

\subsection{Pour un domaine}

Au niveau d'un domaine ou d'un sous-domaine, le cours agrège les réponses de
toutes ses questions avant de calculer les trois mesures. C'est ce qui a été
fait pour le domaine~D au prologue (20 réponses) et au
Chapitre~\ref{chap:decomposition} (sous-domaine par sous-domaine).

\subsection{Quatre cas limites}

\begin{tableaucours}{L{4.6cm}cccL{3.4cm}}{%
    \tableauHeader{Situation} &
    \tableauHeader{$\mu$} & \tableauHeader{$U$} & \tableauHeader{$H$} &
    \tableauHeader{Lecture}
}
Tous répondent \OUI{} & 1 & 0 & 0 & Conformité certaine \\
Tous répondent \NON{} & 0 & 0 & 0 & Non-conformité certaine \\
Tous répondent \NEK{} & 0 & 1 & 0 & \textbf{Ignorance unanime} \\
Réponses également réparties & $1/3$ & $1/3$ & 1,585 & Dispersion maximale \\
\end{tableaucours}
\vspace{-0.8em}
\begin{center}{\small\itshape Tableau~9.1. Les cas limites des trois mesures
(d'après~\cite{minka2026nek})}\end{center}

\begin{boxalert}
\textbf{Violation typique.} Conclure «~certitude~» parce que l'entropie est
nulle.

\smallskip
\textbf{Règle~:} une entropie nulle signifie seulement que tous les
répondants ont donné la même réponse. Si cette réponse est \NEK{}, personne
ne sait rien~: $H = 0$ et $U = 1$. On ne lit jamais $H$ sans $U$ ni $\mu$.
\end{boxalert}

% ============================================================================
\section{Le Quadrant Conformité × Incertitude}
\label{sec:c09-quadrant}
% ============================================================================

L'article A1 propose de lire ensemble la conformité et l'incertitude dans un
quadrant à quatre régions, chacune appelant une action
différente~\cite{minka2026nek}. Le cours retient, par convention, un seuil de
0,5 pour $\mu$ et de 0,20 pour $U$~; chaque mission peut fixer les siens, à
condition de les écrire dans le plan d'audit.

\noindent\begin{minipage}{\linewidth}
\begin{center}
\begin{tikzpicture}[scale=1.0]
\fill[green!18] (3,0) rectangle (6,2);
\fill[red!15] (0,0) rectangle (3,2);
\fill[orange!25] (3,2) rectangle (6,5);
\fill[orange!15] (0,2) rectangle (3,5);
\draw[thick, color=CouleurPrimaire] (0,0) rectangle (6,5);
\draw[dashed, color=CouleurPrimaire] (3,0) -- (3,5);
\draw[dashed, color=CouleurPrimaire] (0,2) -- (6,2);
\node[align=center, font=\footnotesize\bfseries] at (1.5,1.0) {Non-conformité\\ claire};
\node[align=center, font=\footnotesize\bfseries] at (4.5,1.0) {Conformité\\ claire};
\node[align=center, font=\footnotesize\bfseries] at (1.5,3.5) {Non-conformité\\ opaque};
\node[align=center, font=\footnotesize\bfseries] at (4.5,3.5) {Conformité\\ opaque};
\node[font=\scriptsize] at (3,-0.3) {0,5};
\node[font=\scriptsize] at (-0.35,2) {0,20};
\node[font=\small] at (3,-0.8) {Taux de conformité $\mu$};
\node[font=\small, rotate=90] at (-0.9,2.5) {Taux d'incertitude $U$};
\end{tikzpicture}
\end{center}
\vspace{-0.4em}
\begin{center}{\small\itshape Figure~9.1. Le quadrant diagnostique
(d'après~\cite{minka2026nek}~; seuils du cours)}\end{center}
\end{minipage}

\medskip

\begin{tableaucours}{L{3.4cm}L{4.0cm}L{4.8cm}}{%
    \tableauHeader{Région} &
    \tableauHeader{Situation} &
    \tableauHeader{Action}
}
Conformité claire & On s'accorde, avec assurance, sur la conformité &
Validation courante \\
Non-conformité claire & On s'accorde, avec assurance, sur la
non-conformité & Correctifs \\
\textbf{Conformité opaque} & Des \OUI{} nombreux, mais une part importante
de \NEK{} & \textbf{Alerte prioritaire}~: rechercher la preuve avant tout
correctif \\
Non-conformité opaque & Des \NON{} nombreux et une part importante de
\NEK{} & Rechercher la preuve et engager des premiers correctifs \\
\end{tableaucours}
\vspace{-0.8em}
\begin{center}{\small\itshape Tableau~9.2. Les quatre régions et leurs
actions (d'après~\cite{minka2026nek})}\end{center}

La conformité opaque est la région la plus dangereuse~: une grille binaire
la confond avec la conformité claire, alors qu'elle peut cacher une
ignorance répandue. C'est exactement l'erreur du premier rapport de LiZbiZ.

\begin{boxinfo}
\textbf{Deux vocabulaires, une même idée.} L'article \emph{Vers
l'auditologie} distingue la \emph{non-conformité connue}, qui appelle des
correctifs, de la \emph{non-conformité inconnue}, qui appelle d'abord des
investigations~\cite{minka2026auditologie}. La première correspond à la
non-conformité claire~; la seconde aux deux régions opaques du quadrant.
\end{boxinfo}

% ============================================================================
\section{La Règle de Décision du Cours}
\label{sec:c09-regle}
% ============================================================================

\begin{boxprocedure}
\textbf{Conclure sur un domaine.} L'auditeur ne conclut qu'après avoir
consigné, dans le classeur O4, pour chaque domaine~:
\begin{enumerate}
    \item \textbf{$U$ d'abord.} Si $U \geq 0{,}20$, le domaine est opaque~:
          aucune conclusion de conformité n'est possible. L'auditeur
          recherche la preuve, ou conclut \emph{sur l'opacité elle-même}.
    \item \textbf{$\mu$ ensuite.} Si $U < 0{,}20$, le domaine est clair~:
          conforme si $\mu \geq 0{,}5$, non conforme sinon.
    \item \textbf{$H$ pour contrôler.} Une entropie élevée avec $U$ faible
          signale un désaccord entre \OUI{} et \NON{}~: le domaine n'est pas
          homogène (un site conforme, un autre non). On le décompose avant de
          conclure (Chapitre~\ref{chap:decomposition}).
    \item \textbf{Les cotes de l'auditeur prévalent.} Les mesures des
          répondants orientent~; la conclusion repose sur les contrôles cotés
          par l'auditeur, preuve à l'appui (Chapitre~\ref{chap:preuve}).
\end{enumerate}
\end{boxprocedure}

\begin{boxlecture}{Les cinq domaines de LiZbiZ}
\begin{tableaucours}{L{4.0cm}cccL{3.2cm}}{%
    \tableauHeader{Domaine} &
    \tableauHeader{$\mu$} & \tableauHeader{$U$} & \tableauHeader{$H$} &
    \tableauHeader{Région}
}
A. Gouvernance du SI & 0,80 & 0,05 & 0,88 & Conformité claire \\
B. Sauvegardes & 0,67 & 0,08 & 1,19 & Conformité claire, $H$ élevé \\
C. Gestion des stocks & 0,08 & 0,08 & 0,82 & Non-conformité claire \\
D. Contrôle d'accès & 0,65 & 0,25 & 1,24 & \textbf{Conformité opaque} \\
E. Exploitation & 0,40 & 0,30 & 1,57 & Non-conformité opaque \\
\end{tableaucours}

\medskip
\textbf{Analyse} (confiance moyenne).
\begin{itemize}
    \item \emph{A}~: conclusion de conformité.
    \item \emph{B}~: conforme, mais l'entropie signale un désaccord~; la
          décomposition par site montre que l'agence de Garoua ne teste pas
          ses restaurations.
    \item \emph{C}~: non-conformité établie~: correctifs immédiats
          (Chapitre~\ref{chap:controle-interne}).
    \item \emph{D}~: alerte prioritaire~: une conformité apparente cache un
          quart de réponses inconnues~; preuves à rechercher avant toute
          conclusion.
    \item \emph{E}~: investigations et premiers correctifs.
\end{itemize}
\end{boxlecture}

% ============================================================================
\section{Comparer Deux Cycles}
\label{sec:c09-cycles}
% ============================================================================

Un audit ne s'arrête pas au rapport~: le système évolue, et la mission
suivante doit dire s'il s'est amélioré. Le taux de conformité seul ne suffit
pas~: il peut rester stable alors que l'état de connaissance change. L'article
A1 propose la \termecle{divergence de Kullback-Leibler}\index{divergence de
Kullback-Leibler} entre deux cycles~\cite{minka2026nek}~:
\[
    D_{\mathrm{KL}}\!\left(p^{(t+1)} \,\middle\|\, p^{(t)}\right) =
    \sum_{m} p^{(t+1)}(m) \log_2 \frac{p^{(t+1)}(m)}{p^{(t)}(m)}.
\]
Une divergence nulle signifie que la distribution des réponses n'a pas
changé~; une divergence élevée signale un changement qualitatif.

\begin{boxscenario}{Opération LIZBIZ. Le domaine D six mois plus tard}
\textbf{Faits.} En mai~2026, le domaine~D donnait 13~\OUI{}, 2~\NON{},
5~\NEK{}. En novembre~2026, après les corrections, les mêmes questions donnent
18~\OUI{}, 1~\NON{}, 1~\NEK{}.

\medskip
\textbf{Calcul.} $H$ passe d'environ 1,24 à 0,57~bit, ce qui rejoint la valeur
d'environ 0,6~bit que rapporte l'article A0. La divergence entre les deux
cycles vaut
\[
    0{,}90\log_2\frac{0{,}90}{0{,}65} + 0{,}05\log_2\frac{0{,}05}{0{,}10} +
    0{,}05\log_2\frac{0{,}05}{0{,}25} \approx 0{,}26~\text{bit}.
\]

\textbf{Analyse} (confiance moyenne). Le domaine a changé de région~: de la
conformité opaque à la conformité claire. La divergence le confirme.
\end{boxscenario}

\begin{boiteRemarque}{Et si la conformité ne bouge pas~?}
Supposons une question dont la conformité reste à 0,60 sur deux cycles,
mais dont les \NEK{} passent de 0,30 à 0,05 au profit des \NON{}. Le taux de
conformité ne voit rien~; la divergence, elle, vaut environ 0,50~bit.
L'organisation sait désormais ce qu'elle ignorait~: c'est un progrès réel,
même s'il révèle des non-conformités. L'étude de ces trajectoires sur
plusieurs cycles fait l'objet de travaux en préparation.
\end{boiteRemarque}

% ============================================================================
\section{Maturité et Capacité}
\label{sec:c09-maturite}
% ============================================================================

Les directions aiment résumer un domaine par un niveau de 1 à 5. Le cours
propose une conversion du taux de conformité en niveau de maturité, sur
l'échelle à cinq niveaux du modèle CMMI. C'est une \textbf{règle
pédagogique déclarée}, et non une norme.

\begin{tableaucours}{L{2.8cm}cL{3.0cm}L{4.6cm}}{%
    \tableauHeader{$\mu$} &
    \tableauHeader{Niveau} &
    \tableauHeader{Désignation} &
    \tableauHeader{Descripteur}
}
0,90 et plus & 5 & Optimisé & Amélioration continue systématique \\
0,75 à 0,90 & 4 & Quantitatif & Processus mesurés par des indicateurs \\
0,55 à 0,75 & 3 & Défini & Processus standardisés et diffusés \\
0,35 à 0,55 & 2 & Géré & Pratiques documentées et répétables \\
Moins de 0,35 & 1 & Initial & Pratiques ponctuelles, dépendant des
personnes \\
\end{tableaucours}
\vspace{-0.8em}
\begin{center}{\small\itshape Tableau~9.3. Conversion du taux de conformité
en niveau de maturité (convention du cours)}\end{center}

Deux règles encadrent cette conversion~:
\begin{itemize}
    \item \textbf{Pas de niveau pour un domaine opaque.} Si $U \geq 0{,}20$, le
          niveau est \emph{indéterminé}~: un niveau calculé sur des réponses
          inconnues serait trompeur. Chez LiZbiZ, les domaines D et E n'ont
          pas de niveau~; A est au niveau~4, B au niveau~3, C au niveau~1.
    \item \textbf{Ce n'est pas la capacité COBIT.} Le cours GCI
          (MSN-2S3-MSI03) évalue la \emph{capacité} d'un processus COBIT par
          niveaux d'activités, avec l'échelle F/L/P/N et la règle de
          l'ISO/IEC~33000. Le niveau de maturité de ce cours mesure la
          conformité constatée d'un domaine. Les deux échelles répondent à des
          questions différentes et ne se convertissent pas l'une dans l'autre.
\end{itemize}

% ============================================================================
\section{L'Observation en Sept Éléments}
\label{sec:c09-observation}
% ============================================================================

\subsection{Sept éléments}

Chaque non-conformité significative donne lieu à une observation (outil~O8),
de structure invariante.

\begin{tableaucours}{cL{3.0cm}L{8.4cm}}{%
    \tableauHeader{N°} &
    \tableauHeader{Élément} &
    \tableauHeader{Contenu}
}
1 & Norme & Le critère, et le texte ou le référentiel d'où il vient \\
2 & Constatation & Le fait établi, chiffré, avec sa preuve et sa cote \\
3 & Risques & Ce qui peut arriver~: le risque de la fiche d'assertion et sa
conséquence \\
4 & Contradictoire & La réponse de l'audité, datée \\
5 & Analyse & La cause de l'écart et l'interprétation, avec un niveau de
confiance \\
6 & Conclusion & La position finale de l'auditeur \\
7 & Recommandations & Ce qu'il faut faire, en traitant la cause \\
\end{tableaucours}
\vspace{-0.8em}
\begin{center}{\small\itshape Tableau~9.4. Les sept éléments de
l'observation}\end{center}

\subsection{Sept éléments ou cinq~?}

Le cours MSN-2S3-MSI03 présente le constat d'audit en cinq éléments~:
constatation, critère, cause, conséquence, recommandation. Les sept
éléments les développent, sans les contredire.

\begin{tableaucours}{L{5.0cm}L{6.0cm}}{%
    \tableauHeader{Constat en cinq éléments} &
    \tableauHeader{Observation en sept éléments}
}
Critère & 1. Norme \\
Constatation & 2. Constatation \\
Conséquence & 3. Risques \\
· & 4. Contradictoire \\
Cause & 5. Analyse \\
· & 6. Conclusion \\
Recommandation & 7. Recommandations \\
\end{tableaucours}
\vspace{-0.8em}
\begin{center}{\small\itshape Tableau~9.5. Correspondance entre le constat
et l'observation}\end{center}

L'observation ajoute deux éléments que la méthode juge indispensables~: la
trace du contradictoire, qui rend l'observation opposable
(Chapitre~\ref{chap:mission}), et une conclusion explicite, qui sépare la
position de l'auditeur de son analyse.

\begin{boxobjet}{D.2.1 · Observation}
\textbf{1. Norme~:} procédure interne de gestion des accès~: désactivation
du compte d'un agent parti sous 48~heures~; loi n°~2010/012, art.~26.

\textbf{2. Constatation~:} sur 31~départs entre janvier et avril~2026, 19
comptes ont été désactivés sous 48~heures, 7 plus tard, 5 sont encore
actifs (rapprochement exhaustif, papier de travail PT-14). Cote \NON{}.

\textbf{3. Risques~:} usage du compte d'un agent parti pour consulter ou
modifier des données de stock~; manquement à l'obligation de sécuriser le
système.

\textbf{4. Contradictoire~:} le responsable informatique indique, le
2~juillet, que les départs ne lui sont pas notifiés par les ressources
humaines.

\textbf{5. Analyse} (confiance élevée)~: la cause est l'absence de
notification des départs, non la négligence de l'informatique.

\textbf{6. Conclusion~:} le contrôle D.2.1 n'est pas efficace.

\textbf{7. Recommandations~:} notification automatique des départs par les
ressources humaines~; désactivation immédiate des cinq comptes~; revue
trimestrielle des comptes actifs.
\end{boxobjet}

% ============================================================================
\section{Le Rapport et le Suivi}
\label{sec:c09-rapport}
% ============================================================================

\subsection{Le rapport}

Le rapport présente, dans l'ordre~: une synthèse d'une page pour le
commanditaire~; le rappel de la lettre de mission (objectifs, périmètre,
limites)~; la lecture des domaines ($\mu$, $U$, $H$ et région du
quadrant)~; les observations en sept éléments~; la liste hiérarchisée des
recommandations~; en annexe, les réponses écrites des audités.

\begin{boxalert}
\textbf{Violation typique.} Taire, dans la synthèse, les domaines opaques,
pour ne présenter que les non-conformités établies.

\smallskip
\textbf{Règle~:} ce que l'auditeur ne sait pas fait partie de sa conclusion.
Un domaine opaque est signalé comme tel, avec ce qu'il faudrait pour lever
l'incertitude.
\end{boxalert}

\subsection{Le suivi des recommandations}

Une recommandation non suivie est une assurance perdue. Le tableau de suivi
indique, pour chaque recommandation, sa priorité, son responsable, son
échéance, son état (non commencée, en cours, mise en œuvre, vérifiée) et la
pièce qui prouve sa mise en œuvre. La vérification de la mise en œuvre
ouvre le cycle suivant (Chapitre~\ref{chap:mission}).

\begin{boxmetier}{Responsable GRC}
\textbf{Intégrer les conclusions de l'audit}
\begin{enumerate}
    \item Reporter les non-conformités claires dans le registre des risques,
          comme risques dont les contrôles sont inefficaces.
    \item Traiter les domaines opaques comme un risque en soi~: on ne gère
          pas ce que l'on ne connaît pas.
    \item Suivre les recommandations dans le tableau de suivi, et présenter
          leur état au comité d'audit chaque trimestre.
\end{enumerate}
\end{boxmetier}

\separateur

\begin{boxresume}
\textbf{Ce qu'il faut retenir de ce chapitre~:}
\begin{itemize}
    \item \textbf{Trois mesures}~: $\mu$ (part des \OUI{}), $U$ (part des
          \NEK{}), $H$ (dispersion)~; on agrège les réponses pour un domaine.
    \item \textbf{Le quadrant} conformité × incertitude distingue
          conformité et non-conformité, claires ou opaques~; la conformité
          opaque est l'alerte prioritaire.
    \item \textbf{La règle de décision}~: $U$ d'abord, $\mu$ ensuite, $H$ pour
          contrôler~; jamais $H$ seul, car une ignorance unanime donne
          $H = 0$.
    \item \textbf{La divergence de Kullback-Leibler} compare deux cycles et
          voit ce que le taux de conformité ne voit pas.
    \item \textbf{Le niveau de maturité} est une convention déclarée,
          indéterminée pour un domaine opaque, et distincte de la capacité
          COBIT.
    \item \textbf{L'observation en sept éléments} développe le constat en
          cinq~; le rapport dit aussi ce que l'auditeur ne sait pas.
\end{itemize}
\end{boxresume}

\bigskip

\begin{boxexo}[title={\ding{45}\quad Exercice~9.1. Calculer et situer \niveaubase}]
Pour chacune des distributions suivantes (20 réponses), calculez $\mu$, $U$
et $H$, puis situez le domaine dans le quadrant.
\begin{enumerate}[(a)]
    \item 15 \OUI{}, 4 \NON{}, 1 \NEK{}.
    \item 6 \OUI{}, 4 \NON{}, 10 \NEK{}.
    \item 0 \OUI{}, 0 \NON{}, 20 \NEK{}.
    \item 3 \OUI{}, 16 \NON{}, 1 \NEK{}.
\end{enumerate}
\end{boxexo}

\begin{boxexo}[title={\ding{45}\quad Exercice~9.2. Deux cycles à la BANQUE SANAGA \niveaumoyen}]
Sur le domaine « gestion des changements », l'audit interne de la BANQUE
SANAGA obtient 8~\OUI{}, 4~\NON{}, 8~\NEK{} en 2026, puis 12~\OUI{},
6~\NON{}, 2~\NEK{} en 2027.
\begin{enumerate}[(a)]
    \item Calculez $\mu$, $U$ et $H$ pour les deux cycles, et situez le
          domaine dans le quadrant à chaque cycle.
    \item Calculez la divergence de Kullback-Leibler entre les deux cycles.
    \item Le domaine s'est-il amélioré~? Argumentez sans vous limiter au
          taux de conformité.
    \item Quel niveau de maturité attribuer en 2026 et en 2027~?
\end{enumerate}
\end{boxexo}

\begin{boxexo}[title={\ding{45}\quad Exercice~9.3. Le rapport d'AUDIX CONSEIL \niveauexpert}]
À partir de la lecture des cinq domaines de LiZbiZ
(section~\ref{sec:c09-regle}) et des chapitres précédents~:
\begin{enumerate}[(a)]
    \item Rédigez la synthèse d'une page du rapport, en signalant les
          domaines opaques.
    \item Rédigez l'observation en sept éléments relative aux ajustements de
          stock de l'entrepôt de Douala (Chapitre~\ref{chap:fraude}).
    \item Hiérarchisez cinq recommandations en vous appuyant sur le quadrant.
    \item Construisez le tableau de suivi de ces recommandations.
\end{enumerate}
\end{boxexo}

\bigskip

\begin{boxinfo}
\textbf{Transition.} La Partie~\ref{part:methode} a donné la méthode, de la
préparation au suivi. La Partie~\ref{part:normatif} en présente le cadre~:
les normes professionnelles qui l'encadrent, les référentiels qui lui
fournissent ses critères, et le droit camerounais et africain. Le
Chapitre~\ref{chap:normes} ouvre cette partie.
\end{boxinfo}
